\Needspace{6\baselineskip}\SourceLeaf{86}

\[
= X^3\frac{d^2\varphi x}{dt\,du} + 3X^2\frac{dX}{dt}\frac{d\varphi x}{dt}
= \frac{dX^3 \frac{d\varphi x}{dt}}{dt},
\]

\noindent
so folgt

\begin{equation}\tag{5}
\frac{d^3\varphi x}{du^3} = \frac{d^2X^3 \frac{d\varphi x}{dt}}{dt^2}.
\end{equation}

\noindent\hspace*{\parindent}Man findet ebenso

\[
\frac{d^4\varphi x}{du^4} = \frac{d^3X^4 \frac{d\varphi x}{dt}}{dt^3} \quad \text{u.\,s.\,f.}
\]

\Needspace{9\baselineskip}\noindent\hspace*{\parindent}\textls{Allgemein aber} gibt $X^n\dfrac{d\varphi x}{dt}$ nach $u$ differentiirt durch die
Substitutionen (3) und (4)

\begin{equation}\tag{6}
\frac{dX^n \frac{d\varphi x}{dt}}{du} = \frac{dX^{n+1} \frac{d\varphi x}{dt}}{dt}.
\end{equation}

\noindent
Ist also das Gesetz bis auf das $n^{\text{te}}$ Glied richtig oder

\[
\frac{d^n\varphi x}{du^n} = \frac{d^{n-1}X^n \frac{d\varphi x}{dt}}{dt^{n-1}},
\]

\noindent
so wird, nach $u$ differentiirt,

\[
\frac{d^{n+1}\varphi x}{du^{n+1}} = \frac{d\,d^{n-1}X^n\frac{d\varphi x}{dt}}{du . dt^{n-1}}
= \frac{d^{n-1}dX^n\frac{d\varphi x}{dt}}{dt^{n-1}du},
\]

\noindent
[also] nach (6)

\begin{equation}\tag{7}
\frac{d^{n+1}\varphi x}{du^{n+1}} = \frac{d^nX^{n+1} \frac{d\varphi x}{dt}}{dt^n},
\end{equation}

\noindent
d.\,i.\ [die Regel gilt] auch für das nächstfolgende Glied.

\noindent\hspace*{\parindent}Aus den Gleichungen (1), (4), (5) folgen also die obigen Verwandlungen für die drei ersten
Glieder, und aus (6) ihre Richtigkeit für alle folgenden.

\medskip

\Needspace{10\baselineskip}\begin{center}
\textbf{BEMERKUNG ZUM AUFSATZE »NEUER BEWEIS DES LAGRANGISCHEN LEHRSATZES«.}
\end{center}

\noindent\hspace*{\parindent}Zufolge einer Notiz in dem S. 20 erwähnten Tagebuche hat \textsc{Gauss} am 27. December 1796 einen
Beweis des Lehrsatzes von \textsc{Lagrange} gefunden. Der vorstehend abgedruckte Aufsatz ist zwar
ohne jede Datumangabe auf einem einzelnen Zettel aufgezeichnet; indessen ist nicht zu
bezweifeln, dass es sich hier um die Entwicklung handelt, welche \textsc{Gauss} Ende 1796 auffand.
Die Veröffentlichung wurde seiner Zeit nur durch verschiedene Zufälligkeiten verhindert
(Vgl.\ \textsc{Sartorius von Waltershausen}, Gauss zum Gedächtniss, S. 22).

\begin{flushright}
\textsc{Fricke}.
\end{flushright}
